\usepackage[utf8]{inputenc} % Kodierung
\usepackage[T1]{fontenc} % Explizite Nennung des Fonts
% Use Times-like font (text + math)
\usepackage{newtxtext,newtxmath}
\usepackage[ngerman]{babel} % Sprache (Deutsch)
\usepackage{graphicx} % immer benötigt für das Einbinden von Graphiken
\usepackage{blindtext} % Wenn man das Layout prüfen will, kann hier mit \blindtext Text eingfügt werden.
%\usepackage{parskip} % Für den Abstand zwischen 2 Absätzen.
%\setlength{\parskip}{7pt plus20pt minus10pt} % Genaue Einstellung von parskip
\usepackage{subcaption}
\usepackage{easy-todo} % Mit \todo{} Todos einfügen
\usepackage{csquotes} % Für ordentlichen Anführungszeichen
\usepackage[iso, ngerman]{isodate} % Deutsche Datumsformatierung
\usepackage[style=apa, 
backend=biber, 
bibencoding=utf8, 
language=ngerman, 
mincitenames=1, 
maxcitenames=2]{biblatex}
\addbibresource{literatur/bibliography.bib} % Einbinden der Literatur.
\AtEveryBibitem{\clearfield{annotation}} % Unterdrückt "Article." aus den annote-Feldern der .bib-Datei
\usepackage{amsmath}

\DefineBibliographyStrings{ngerman}{%
  urlseen = {Abgerufen am},
  % APA verlangt "et al." statt der deutschen Abkuerzung "u. a."
  andothers = {et\addabbrvspace al\adddot},
}

\DeclareLanguageMapping{ngerman}{ngerman-apa} % Bibliographie-Sprache fuer APA
\usepackage{float} % Grafik bleibt dort wo man sie im Code platziert hat

\usepackage[activate={true,nocompatibility},
	final,
	tracking=true,
	kerning=true,
	expansion=true,
	spacing=true,
	factor=1050,
	stretch=25,
	shrink=10]{microtype} % Für die Feineinstellung der Zeichensetzung.
\usepackage{booktabs}
\usepackage{appendix}
\usepackage[rflt]{floatflt}
\usepackage{fancyvrb}
\usepackage[hidelinks]{hyperref} % Klickbare aber nicht markierte Links im PDF
\usepackage{setspace}
\usepackage{fancyhdr} % Für schönere Kopf-/Fußzeilen und Fußnoten.
\usepackage[left=3cm, right=2cm, top=2.5cm, bottom=2cm]{geometry} % Seitenränder
\usepackage{pbox}
\usepackage{tabulary}
\usepackage{outlines}
\usepackage{matlab-prettifier}
\usepackage{listings}
\usepackage{color} %red, green, blue, yellow, cyan, magenta, black, white

% Für die Abbildung
\usepackage{tikz}
\usepackage{pgfplots}
\pgfplotsset{width=0.7*\textwidth,compat=1.9}

% Für Tabellen
\usepackage{csvsimple}
\usepackage{longtable}
\usepackage{rotating}
\usepackage{pdflscape} % Querformat-Seiten (inkl. gedrehter Überschrift)
\usepackage{multirow}

\definecolor{mygreen}{RGB}{28,172,0} % color values Red, Green, Blue
\definecolor{mylilas}{RGB}{170,55,241}

%Kein Schusterjunge/Hurenkind im Text
\widowpenalty10000
\clubpenalty10000
\usepackage[justification=centering]{caption}

\usepackage{titlesec}

% --- Formatierung für \section ---

% Aussehen der Überschrift
\titleformat{\section}
  {\normalfont\LARGE\bfseries} % Hier die Größe anpassen (z. B. \Large, \LARGE, \huge)
  {\thesection}                % Die Nummerierung (z.B. "1")
  {1em}                        % Abstand zwischen Nummerierung und Text
  {}                           % Zusätzlicher Code vor der Überschrift (hier leer)

% Abstände der Überschrift
\titlespacing*{\section}
  {0pt}      % Einzug von links
  {18pt}     % Abstand nach oben (vor der Überschrift)
  {18pt}     % Abstand nach unten (nach der Überschrift)
  \titlespacing*{\subsection}
  {0pt}      % Einzug von links
  {12pt}     % Abstand nach oben (vor der Überschrift)
  {12pt}     % Abstand nach unten (nach der Überschrift)
% Aussehen der \paragraph-Ueberschrift (eigene Zeile statt Fliesstext)
\titleformat{\paragraph}[hang]
  {\normalfont\normalsize\bfseries}
  {\theparagraph}
  {1em}
  {}
\titlespacing*{\paragraph}
  {0pt}
  {10pt}
  {6pt}

% Abstand ueber und unter abgesetzten Formeln.
% Standard sind 12pt ueber und unter dem Block, waehrend einzelne Gleichungen nach
% einer kurzen Zeile mit \abovedisplayshortskip = 0pt gesetzt werden. Dadurch stehen
% align-Bloecke deutlich lockerer als \begin{equation}. Die Werte unten gleichen das an.
% \normalsize setzt die Laengen bei jedem Schriftgroessenwechsel auf die Klassenwerte
% zurueck, deshalb werden sie dort angehaengt.
\makeatletter
\newcommand{\FormelAbstaende}{%
  \setlength{\abovedisplayskip}{5pt plus 2pt minus 2pt}%
  \setlength{\belowdisplayskip}{8pt plus 2pt minus 3pt}%
  \setlength{\abovedisplayshortskip}{6pt plus 2pt minus 1pt}%
  \setlength{\belowdisplayshortskip}{8pt plus 2pt minus 3pt}%
}
\g@addto@macro\normalsize{\FormelAbstaende}
\makeatother
\AtBeginDocument{\FormelAbstaende}
